% ============================================================
% 文章式课程报告模板（uestc-create · latex-article）
% 体例：ctexart 文章式（\section 制），12pt、1.5 倍行距、
%       GB/T 7714-2015 著者-出版年引用（(作者, 年份)）。
% 适用：普通课程论文/小论文/课程报告；正式大报告/结题建议用
%       上一级 ../latex/（thesis-uestc 仿学位论文章节制）。
% 编译：xelatex 两遍；改动 .bib 后 xelatex → bibtex → xelatex ×2。
% 说明：封面与正文均为占位文字（张三/XX课程/20XX），按课程要求替换。
% ============================================================
\documentclass[UTF8,a4paper,12pt]{ctexart}

\usepackage{mathtools}
\usepackage{amsmath}
\usepackage{amssymb}
\usepackage{bm}
\usepackage{graphicx}
\usepackage{booktabs}
\usepackage{multirow}
\usepackage{enumerate}
\usepackage{titlesec}
\usepackage{titletoc}
\usepackage{geometry}
\usepackage{fancyhdr}
\usepackage{setspace}
\usepackage[colorlinks,linkcolor=blue,anchorcolor=blue,citecolor=green]{hyperref}

\geometry{left=2.5cm, right=2.5cm, top=2.8cm, bottom=2.8cm}
\graphicspath{{images/}}
\onehalfspacing

\newcommand{\HRule}{\rule{\linewidth}{0.5mm}}

% 目录样式：点线引导 + 页码右对齐
\titlecontents{section}
              [0em]
              {\bf \small\vspace{2pt}}%
              {\contentslabel{2.5em}}%
              {}%
              {\titlerule*[0.5pc]{$\cdot$}\contentspage}%
\titlecontents{subsection}
              [2.5em]
              {\small}%
              {\contentslabel{2.5em}}%
              {}%
              {\titlerule*[0.5pc]{$\cdot$}\contentspage}%

\begin{document}

% ============================================================
% 封面（占位字段：学院/课程/题目/作者/日期，请全部替换）
% ============================================================
\begin{titlepage}
\begin{center}
\includegraphics[width=\textwidth]{UESTC-Logo}\\[1cm]
\textsc{\LARGE XX学院~$\bullet$~XX课程}\\[1.5cm]
\HRule \\[0.4cm]
{\huge \bfseries 课程论文题目\\[0.4cm]}
\HRule \\[1.5cm]
\noindent
\emph{作者}:\quad {\bf 张三}\footnotemark[1]\\[0.5cm]
\emph{题目}:\quad 论文副标题或英文题目（占位）\\[0.5cm]
\emph{日期}:\quad 20XX 年 X 月
\end{center}
\footnotetext[1]{电子科技大学XX学院XX专业}
\end{titlepage}

% ============================================================
% 目录
% ============================================================
{
\small
\setstretch{1.15}
\tableofcontents
}
\newpage

% ============================================================
% 正文（各节均为占位文字，按"观点句 → 证据 → 小结"展开；
%       每段 1–2 处引用为度，编号与文末条目一一对应）
% ============================================================
\section{引言}

占位：背景（带事实或数字）→ 已有工作 → 缺口或问题 → 本文做什么、怎么做，结尾给一句路线声明。引用示例\cite{zhangsan2024}。

\section{方法}

占位：每小节回答一个明确的小问题。公式示例：
\begin{equation}
  L = \frac{1}{N}\sum_{i=1}^{N} \ell\!\left(f(x_i;\theta),\, y_i\right)
\end{equation}
式中 $L$ 为损失（占位），$N$ 为样本数（占位）。

\begin{figure}[!ht]
    \centering
    % \includegraphics[width=0.85\textwidth]{example.png}
    % 图片放 images/，正文只写文件名；删掉本注释并换成真实图片
    \fbox{图片占位（images/example.png）}
    \caption{图题占位（图题在图下，全文连续编号）}
    \label{fig:example}
\end{figure}

\section{实验与分析}

占位：数字先行，每图每表必有一段解读。表格示例：
\begin{table}[!ht]
\centering
\caption{表题占位（表题在表上）}
\label{tab:example}
\begin{tabular}{lcc}
\toprule
方法 & 指标A & 指标B \\
\midrule
基线 & 0.00 & 0.00 \\
本文 & 0.00 & 0.00 \\
\bottomrule
\end{tabular}
\end{table}

\section{总结}

占位：做出了什么、意味着什么、还剩什么；不复述方法。

% ============================================================
% 参考文献（GB/T 7714-2015 著者-出版年；条目见 Reference.bib）
% ============================================================
\bibliographystyle{gbt-7714-2015-author-year}
\bibliography{Reference}

\end{document}
